\phantomsection
\addcontentsline{toc}{section}{About Part II}
\section*{About Part II}

Part II is where the degree stops being a degree and becomes a menu. There are around
forty courses on offer and no obligation to take any particular one, so no two students
in the year sit the same set. The twenty-one chapters here are the ones I took, and
they are not a syllabus so much as a set of choices; another student's Part II could
overlap with this one in three courses and be no less complete.

The change in character is real, and it is not only about choice. Part IB abstracted
the objects of Part IA; Part II mostly stops abstracting and starts going deep. A
typical course here takes one idea and follows it for twenty-four lectures until it
arrives somewhere that would have been unimaginable two years earlier -- Galois Theory
begins with field extensions and ends by proving that the general quintic is not
soluble by radicals; Algebraic Geometry begins with the zero set of a polynomial and
ends at the Riemann--Roch theorem; Algebraic Topology begins with loops in a space and
ends by proving that $\R^m \cong \R^n$ forces $m = n$. The courses also stop being
polite to one another. In Part IB every course could reasonably assume every other; in
Part II a lecturer has to assume a specific list, and the list is short.

The algebraic strand is the one that most obviously pays off. Groups, Rings and Modules
splits into four: Galois Theory, which turns questions about polynomials into questions
about finite groups; Representation Theory, which turns questions about finite groups
into questions about matrices, and gets character tables and Burnside's theorem out of
it; Number Fields, which repairs unique factorisation in rings of algebraic integers by
replacing elements with ideals; and Algebraic Geometry, which builds a dictionary
between ideals and varieties and then uses it relentlessly. Number Theory sits slightly
below these, elementary in its methods and enormously useful in its results, and feeds
both Number Fields and Coding and Cryptography.

The analytic strand pulls in the opposite direction, towards infinite dimensions.
Linear Analysis is what happens when linear algebra is done on a Banach space, and its
three great theorems -- Hahn--Banach, Baire category and the consequences of
completeness -- are the tools of everything after it. Probability and Measure rebuilds
the whole of Part IA Probability on a measure-theoretic foundation, which is what
finally makes the strong law of large numbers and the central limit theorem provable
rather than plausible. Analysis of Functions then combines the two, developing $L^p$
spaces, weak topologies, distributions and Sobolev spaces, and ends by solving elliptic
partial differential equations with them. Topics in Analysis is the outlier of the
group and a delight: a collection of essentially unrelated beautiful results -- fixed
point theorems, Runge's theorem, Gaussian quadrature, the irrationality of $e$ and
$\pi$ -- held together by nothing but taste.

Geometry in Part II is where the various strands meet. Differential Geometry retells IB
Geometry with more machinery and more nerve, arriving at Gauss--Bonnet for the second
time in this book. Riemann Surfaces is the most cross-bred course of the year, needing
complex analysis for its local theory, covering spaces from Algebraic Topology for its
global theory, and pointing directly at Algebraic Geometry. Algebraic Topology itself
develops the fundamental group and then simplicial homology, and is the course that
best explains why anybody ever wanted a functor.

The mathematical physics courses continue the IB story. Principles of Quantum Mechanics
recasts IB Quantum Mechanics in Dirac's language, where the harmonic oscillator, angular
momentum and perturbation theory become algebra rather than differential equations, and
it in turn supports Statistical Physics, which is the one course in the book that
genuinely derives thermodynamics rather than assuming it. Integrable Systems occupies
strange and wonderful territory between the applied and the pure, using the inverse
scattering transform to solve nonlinear PDEs exactly. Quantum Information and
Computation needs almost nothing but linear algebra and a little quantum mechanics, and
gets to Shor's and Grover's algorithms remarkably fast.

The remaining courses are the genuinely independent ones. Logic and Set Theory builds
ordinals, cardinals and the axiom of choice and then proves G\"odel's incompleteness
theorem; Automata and Formal Languages does the computability half of the same story;
Graph Theory is combinatorics with almost no prerequisites and some of the best proofs
in the book; Coding and Cryptography turns Shannon's theorems into practical codes;
Stochastic Financial Models takes the probability strand into continuous time and
derives Black--Scholes. Any of these five can be read cold, given Part IA and a little
patience.

If Part IA rewards being read in order and Part IB rewards being read in strands, Part
II rewards being read out of curiosity. Pick the chapter whose first page interests
you, and use the dependency graphs in the front matter to work out what you owe.
